\section{Experiments}
\label{sec:experiments}

\paragraph{Setup.}
We use ViT-Small/16~\cite{dosovitskiy2021image} pretrained on ImageNet-21k
(21.7M params) via \texttt{timm}~\cite{wightman2021timm}.
LoRA adapters ($r{=}4$) are applied to $W_Q, W_K, W_V, W_O$ in all 12 attention
blocks (307K trainable, 1.41\% of total); the classifier head is fully
fine-tuned.
We evaluate on \textbf{FGVC-Aircraft}~\cite{maji2013fgvc} (100 variants;
extreme similarity), \textbf{CUB-200-2011}~\cite{wah2011cub} (200 birds;
moderately separable), and \textbf{Stanford
Dogs}~\cite{khosla2011stanforddogs} (120 breeds).
All models train for 50 epochs with AdamW~\cite{loshchilov2019adamw}
($\eta{=}5{\times}10^{-4}$, cosine decay); EWL defaults are $\alpha{=}0.9$,
$\tau{=}1.0$, warmup $K{=}48$.
Results are mean\,$\pm$\,std over three seeds $\{11,22,33\}$.
We compare \textbf{SFT} (LoRA baseline), \textbf{EWL} (full method), and
\textbf{EWL (no proxy)} ($g\equiv 1$).
Full protocols are in Appendix~\ref{app:implementation}.

\paragraph{Clean accuracy and proxy ablation.}
On clean data, EWL with the proxy matches SFT within standard deviation on
all three benchmarks ($\Delta\!=\!+0.05, +0.16, -0.01$); we do not claim
clean-data gains.
The decisive result is the no-proxy ablation: removing the gradient proxy
while retaining the progress signal causes a $\mathbf{-6.76\%}$ drop on
Aircraft, with smaller but consistently negative drops on CUB-200 and Dogs.
Aircraft's 100 visually similar classes have densely packed decision boundaries;
indiscriminate early reweighting, before the EMA has stabilised, amplifies
unstable gradient directions and the model cannot recover.
Full numbers are in Appendix~\ref{app:curves} (Table~\ref{tab:main_results});
training curves there confirm the gap is a training instability rather than a
final-epoch artefact.

\paragraph{Noise robustness.}
We inject uniform random label noise into the training set at rates of 10--40\%
and report validation accuracy at the earliest epoch where EWL matches or
exceeds the SFT baseline (Table~\ref{tab:noise_results},
Figure~\ref{fig:noise_delta}).
On Aircraft, EWL improves over SFT by $+4.26\%$ to $+6.81\%$ across all noise
levels and does so consistently at the final epoch.
On CUB-200, EWL reaches a $+1.92\%$ / $+2.82\%$ advantage over SFT at
20\%/40\% noise respectively, but reaches this advantage several epochs
earlier than SFT, suggesting EWL converges faster under label noise.
On Stanford Dogs, EWL matches SFT within noise ($+0.04\%$, $+0.12\%$)
at the same matched epoch.
The magnitude of the Aircraft gain --- and its persistence through final
training --- is the central empirical finding; the mechanism is analysed in
\S\ref{sec:discussion}.
Full noise results at every level are in Appendix~\ref{app:noise}.

\begin{table}[t]
  \centering
  \small
  \caption{
    Top-1 validation accuracy (\%) under label noise, mean\,$\pm$\,std over
    3 seeds. For Aircraft we report final-epoch test accuracy; for CUB-200
    and Dogs we report validation accuracy at the earliest epoch where EWL
    matches or exceeds the SFT baseline.
    Aircraft shows a large, monotonically increasing EWL advantage;
    CUB-200 shows a smaller positive advantage realised earlier in training;
    Dogs matches SFT within noise.
  }
  \label{tab:noise_results}
  \setlength{\tabcolsep}{4pt}
  \begin{tabular}{llccc}
    \toprule
    \textbf{Dataset} & \textbf{Noise} & \textbf{SFT} & \textbf{EWL} & $\boldsymbol{\Delta}$ \\
    \midrule
    FGVC-Aircraft & 10\% & $51.90{\pm}0.65$ & $56.16{\pm}0.43$ & $\mathbf{+4.26}$ \\
    FGVC-Aircraft & 20\% & $46.21{\pm}0.98$ & $52.12{\pm}0.55$ & $\mathbf{+5.90}$ \\
    FGVC-Aircraft & 30\% & $40.26{\pm}0.64$ & $47.07{\pm}0.66$ & $\mathbf{+6.81}$ \\
    FGVC-Aircraft & 40\% & $34.64{\pm}1.03$ & $40.89{\pm}1.68$ & $\mathbf{+6.25}$ \\
    \midrule
    CUB-200-2011  & 20\% & $73.79{\pm}0.15$ & $75.71{\pm}0.75$ & $\mathbf{+1.92}$ \\
    CUB-200-2011  & 40\% & $66.69{\pm}1.40$ & $69.51{\pm}0.38$ & $\mathbf{+2.82}$ \\
    \midrule
    Stanford Dogs & 20\% & $85.67{\pm}0.21$ & $85.71{\pm}0.31$ & $+0.04$ \\
    Stanford Dogs & 40\% & $83.26{\pm}0.31$ & $83.38{\pm}0.45$ & $+0.12$ \\
    \bottomrule
  \end{tabular}
\end{table}

\begin{figure}[t]
  \centering
  \includegraphics[width=0.95\linewidth]{figures/fig_noise_delta_lines_all_datasets.pdf}
  \caption{
    Accuracy delta $\Delta(\text{EWL}-\text{SFT})$ vs.\ label noise rate.
    Aircraft rises to $+6.81\%$ at 30\% noise; CUB-200 shows a smaller
    $+1.92$ to $+2.82\%$ advantage at the matched epoch; Stanford Dogs
    matches SFT within noise.
    The scale of the Aircraft gain is the central empirical finding of this
    work.
  }
  \label{fig:noise_delta}
\end{figure}

\paragraph{Ablations and overhead.}
We sweep LoRA rank $r \in \{2,4,8,16\}$, softmax temperature
$\tau \in \{0.5,\dots,3.0\}$, EMA decay $\alpha \in \{0.5,\dots,1.0\}$, and
class imbalance $\text{IF}\in\{1,2,5,10,20\}$ on Aircraft (all tables and
figures in Appendix~\ref{app:hparam}--\ref{app:imbalance}).
Summarising: accuracy increases monotonically with rank ($57.40$ at $r{=}2$ to
$62.08$ at $r{=}16$); $\tau{=}0.5$ over-concentrates and hurts, while
$\tau\!\geq\!1.0$ is stable within $0.2\%$; EWL provides no meaningful benefit
under class imbalance.
Step-time overhead is $5.9\times$ SFT (stable $\sim$265\,ms vs.\ $\sim$45\,ms)
with \textbf{zero} GPU memory overhead: both conditions peak at 9.46\,GB.
The cost comes from CPU-side per-sample operations (EMA update, signal,
softmax), serialised by CPU$\leftrightarrow$GPU synchronisation rather than
floating-point throughput; vectorising these onto GPU tensors is natural
future work.
